
\documentclass{article}
\usepackage{colortbl}
\usepackage{makecell}
\usepackage{multirow}
\usepackage{supertabular}

\begin{document}

\newcounter{utterance}

\centering \large Interaction Transcript for game `cladder', experiment `full\_v1.5\_default', episode 5313 with qwen3.5{-}4b.
\vspace{24pt}

{ \footnotesize  \setcounter{utterance}{1}
\setlength{\tabcolsep}{0pt}
\begin{supertabular}{c@{$\;$}|p{.15\linewidth}@{}p{.15\linewidth}p{.15\linewidth}p{.15\linewidth}p{.15\linewidth}p{.15\linewidth}}
    \# & \multicolumn{2}{c}{Player} && \multicolumn{2}{c}{Game Master} \\
    \hline

    \theutterance \stepcounter{utterance}  
    & & & \multicolumn{4}{p{0.6\linewidth}}{
        \cellcolor[rgb]{0.9,0.9,0.9}{
            \makecell[{{p{\linewidth}}}]{
                \texttt{\tiny{[P1$\langle$GM]}}
                \texttt{You are an expert in causal inference. The following question is not a typical commonsense query, but rather a meticulously designed question created by a professor specializing in causal inference, intended to assess the students' mastery of the course content.} \\
\texttt{Imagine a self{-}contained, hypothetical world with only the following conditions, and without any unmentioned factors or causal relationships: Liking spicy food has a direct effect on occupation and salary. Skill has a direct effect on occupation and salary. Occupation has a direct effect on salary. Skill is unobserved. The overall probability of liking spicy food is 2\%. The probability of not liking spicy food and high salary is 45\%. The probability of liking spicy food and high salary is 1\%. Is the chance of high salary smaller when observing liking spicy food?} \\
\texttt{Start your answer with 'Yes' or 'No', followed by additional reasoning or evidence to support your explanation.} \\
            }
        }
    }
     \\ \\

    \theutterance \stepcounter{utterance}  
    & \multicolumn{4}{p{0.6\linewidth}}{
        \cellcolor[rgb]{0.561,0.737,0.561}{
            \makecell[{{p{\linewidth}}}]{
                \texttt{\tiny{[P1$\rangle$GM]}}
                \texttt{Yes} \\
\\ 
\texttt{**Reasoning:**} \\
\\ 
\texttt{To determine whether the chance of high salary is smaller when observing that someone likes spicy food, we must compare the conditional probability of high salary given spicy food (\$P(\textbackslash{}text\{High Salary\} | \textbackslash{}text\{Spicy\})\$) with the conditional probability of high salary given no spicy food (\$P(\textbackslash{}text\{High Salary\} | \textbackslash{}text\{Not Spicy\})\$).} \\
\\ 
\texttt{Let \$S\$ denote the event of liking spicy food and \$H\$ denote the event of having a high salary.} \\
\texttt{We are given:} \\
\texttt{1.  \$P(S) = 0.02\$} \\
\texttt{2.  \$P(\textbackslash{}neg S \textbackslash{}cap H) = 0.45\$} \\
\texttt{3.  \$P(S \textbackslash{}cap H) = 0.01\$} \\
\\ 
\texttt{First, let's calculate the marginal probability of having a high salary, \$P(H)\$, using the law of total probability:} \\
\texttt{\$\$P(H) = P(S \textbackslash{}cap H) + P(\textbackslash{}neg S \textbackslash{}cap H)\$\$} \\
\texttt{\$\$P(H) = 0.01 + 0.45 = 0.46\$\$} \\
\\ 
\texttt{Next, we calculate the conditional probability of high salary given that the person likes spicy food:} \\
\texttt{\$\$P(H | S) = \textbackslash{}frac\{P(S \textbackslash{}cap H)\}\{P(S)\} = \textbackslash{}frac\{0.01\}\{0.02\} = 0.5\$\$} \\
\\ 
\texttt{Then, we calculate the conditional probability of high salary given that the person does *not* like spicy food:} \\
\texttt{\$\$P(H | \textbackslash{}neg S) = \textbackslash{}frac\{P(\textbackslash{}neg S \textbackslash{}cap H)\}\{P(\textbackslash{}neg S)\}\$\$} \\
\\ 
\texttt{To find \$P(\textbackslash{}neg S)\$, we use the complement rule:} \\
\texttt{\$\$P(\textbackslash{}neg S) = 1 {-} P(S) = 1 {-} 0.02 = 0.98\$\$} \\
\\ 
\texttt{Now, substitute the values into the conditional probability formula:} \\
\texttt{\$\$P(H | \textbackslash{}neg S) = \textbackslash{}frac\{0.45\}\{0.98\} \textbackslash{}approx 0.4592\$\$} \\
\\ 
\texttt{**Comparison:**} \\
\texttt{*   \$P(H | S) = 0.50\$} \\
\texttt{*   \$P(H | \textbackslash{}neg S) \textbackslash{}approx 0.459\$} \\
\\ 
\texttt{Since \$0.50 > 0.459\$, the probability of having a high salary is actually **higher** (not smaller) when observing that someone likes spicy food. The observation of liking spicy food is positively associated with high salary in this specific hypothetical world, despite the unobserved confounder "Skill" which likely drives both the preference for spicy food and the occupation/salary.} \\
\\ 
\texttt{Therefore, the answer to the question "Is the chance of high salary smaller when observing liking spicy food?" is **No**.} \\
            }
        }
    }
    & & \\ \\

    \theutterance \stepcounter{utterance}  
    & & & \multicolumn{2}{p{0.3\linewidth}}{
        \cellcolor[rgb]{0.9,0.9,0.9}{
            \makecell[{{p{\linewidth}}}]{
                \texttt{\tiny{[GM$|$GM]}}
                \texttt{Yes} \\
            }
        }
    }
    & & \\ \\

    \theutterance \stepcounter{utterance}  
    & & & \multicolumn{2}{p{0.3\linewidth}}{
        \cellcolor[rgb]{0.9,0.9,0.9}{
            \makecell[{{p{\linewidth}}}]{
                \texttt{\tiny{[GM$|$GM]}}
                \texttt{yes} \\
            }
        }
    }
    & & \\ \\

    \theutterance \stepcounter{utterance}  
    & & & \multicolumn{2}{p{0.3\linewidth}}{
        \cellcolor[rgb]{0.9,0.9,0.9}{
            \makecell[{{p{\linewidth}}}]{
                \texttt{\tiny{[GM$|$GM]}}
                \texttt{game\_result = WIN} \\
            }
        }
    }
    & & \\ \\

\end{supertabular}
}

\end{document}
